% 本文件由 LaTeX 排版 Agent 根据有效 translation.json 自动生成。
% author=Councils; work_id=3805; title=安提约基雅奉献节会议（公元341年）

\chapter{安提约基雅奉献节会议（公元341年）}

% structure: p0001 source=h1 target=omitted reason=page-title-already-rendered-by-parent

% structure: p0002 source=h2 target=section reason=relative-heading-rank; base=h2

\section{历史引言}

关于采纳了随后被纳入普世教会法典的诸法条的安提约基雅会议，我们知道其确切日期。这一日期由以下事实确定：会议是在安提约基雅的大教堂——被称为\Quote{金殿}——的奉献典礼期间举行的，该教堂由其父君士坦丁大帝始建，在君士坦提乌斯时代完工。因此，该会议一直被称为安提约基雅\OriginalTerm{奉献节}{in\ Encæniis}会议，即\OriginalTerm{奉献典礼}{in\ Dedicatione}期间，于公元341年夏季举行。九十七位主教聚集一堂，其中许多人敌视\Person{圣亚他那修}，他们自称\OriginalTerm{优西比派}{Eusebians}，全是东方主教，且大部分隶属于安提约基雅宗主教区。没有一位西方或拉丁主教出席，教宗\Person{犹利}也未以任何方式代表出席。这一事实给了史学家\Person{苏格拉底}机会作出如下陈述（围绕此陈述曾爆发激烈论战）：\Quote{一条教会法令规定，各教会不得违背罗马主教的意见作出决议。}\par

然而，尽管这一切都清楚明了，但没有任何一个会议给史学家和神学家带来如此之多的困难。无人能否认，时人\Person{普瓦捷的圣依拉略}称其为\OriginalTerm{圣者议会}{Synodus\ Sanctorum}；其两条法条在加采东被宣读为\Quote{圣教父法条}；教宗\Person{若望二世}、\Person{匝加利亚}和\Person{良四世}都曾批准这些法条，并将其归诸\Quote{圣教父}。然而，这个会议却制定了与尼西亚会议相抗衡的信经，并且据说有些法条是为了谴责\Person{亚他那修}而采纳的。\par

人们曾尝试各种办法以摆脱这些困难。\par

有人提出，实际上有两个安提约基雅会议，一个正统，制定了法条；另一个是异端。\par

\Person{埃马努埃尔·谢尔斯特拉滕神父（S.J.）}完善了这一理论。他设想，在正统主教离开后，\OriginalTerm{优西比派}{Eusebians}留在安提约基雅，随后通过了反对\Person{亚他那修}的决议，并宣称会议仍在进行。此理论为\Person{帕吉}、\Person{雷米·切利耶}、\Person{瓦尔希}所采纳，并在一定程度上为\Person{施勒克}等人所采纳。但\Person{蒂尔蒙}对此持异议，主张据\Person{苏格拉底}所述，\Person{亚他那修}的罢黜在先，法条的采纳在后。但\Person{蒂尔蒙}似乎在此点上误解了\Person{苏格拉底}，因此这一异议便不成立。但另一异议仍然存在，即\Person{苏格拉底}和\Person{索措门}都说信经是在\Person{亚他那修}被罢黜\Emph{之后}才起草的，\Quote{然而}（\Person{赫费勒}指出，第二卷，第63页），\Quote{\Person{圣依拉略}说这些信经出自\InnerQuote{圣者议会}。}\par

\Person{谢尔斯特拉滕}的假设不能令人满意，博学的\Person{巴莱里尼}兄弟在其\WorkTitle{圣大良全集}的附录中提出了另一理论，\Person{曼西}在其\Quote{对\Person{亚历山大·纳塔利斯}《\WorkTitle{教会史}》的注释}中同意此理论。他们认为，这些法条根本不出自\OriginalTerm{奉献节}{in\ Encæniis}会议，而是出自早先于332年举行的另一个会议；但\Person{赫费勒}全然拒绝这一假设，理由如下：首先且最重要的是，它没有外部证据支持；其次，内部证据极不令人满意。然而，即使这二十五条法条是由332年的一次安提约基雅会议所采纳，真正的困难仍不会消除，因为\Person{苏格拉底}论及那次会议时说，在那里\Quote{反对尼西亚信仰的人}也能选出其候选人，以填补被放逐的主教\Person{优斯塔修斯}之位！\par

\Person{赫费勒}似乎在说出整个难题的真正解法时说道：\Quote{诚然，\Person{亚他那修}将\OriginalTerm{优西比派}{Eusebians}等同于\OriginalTerm{亚略派}{Arians}，而我们视他们至少为\OriginalTerm{半亚略派}{Semi\=|arians}；但在当时，在他们作了正统的信仰宣认，并多次声明反对在尼西亚被谴责的异端之后，他们被大多数人视为合法的主教，并且完全正统和圣洁的人士可以毫不犹豫地在会议中与他们联合。}\par

教宗\Person{犹利}在指责其行为的同时，称那罢黜了\Person{亚他那修}的优西比派会议为\Quote{亲爱的弟兄们}，并邀请他们参加一个共同会议，以调查针对这位圣人的指控。鉴于这一切，我们完全可以相信，正统派和优西比派都聚集在一起参加了皇帝新教堂的祝圣典礼，而后来整个教会根据这些法条的内在价值赋予了它们相应的地位，而不考虑那些起草和投票支持它们的人的动机或神学观点的细微差别。\par

% structure: p0011 source=h2 target=section reason=relative-heading-rank; base=h2

\section{会议书信}

从\Place{科勒叙利亚}、\Place{腓尼基}、\Place{巴勒斯坦}、\Place{阿拉伯}、\Place{美索不达米亚}、\Place{基利家}和\Place{以扫利亚}各省聚集在安提约基雅的圣洁而极其和平的会议，致各教省同心合意且圣洁的同工：愿在主内安康。\par

我等主救主\Person{耶稣基督}的恩宠与真理眷顾了安提约基雅人的圣教会，以同心同德与和平之圣神将其联合，同样也改善了其他许多事物；在这一切之中，藉着圣洁而赐予和平的圣神的助佑，这种改善得以实现。因此，我们这些来自不同省份、在安提约基雅聚集的主教们，经过多次审查和考察后一致同意的事，现在告知诸位；我们信赖基督的恩宠和和平的圣神，相信你们也会同意我们，并尽己所能与我们站在一起，以祈祷与我们一同奋斗，更与你们联合，跟随圣神，在我们的定义上同心，制定与我们相同的事；你们在圣神所赐的和谐中，将所决定的事盖上印并确认。\par

现将所制定的教会法条附录于后。\par

% structure: p0015 source=h2 target=section reason=relative-heading-rank; base=h2

\section{法条}

% structure: p0016 source=h3 target=subsection reason=relative-heading-rank; base=h2

\subsection{第一条}

\Emph{凡}胆敢废止在敬爱天主的虔诚皇帝\Person{君士坦丁}面前、于尼西亚召开的神圣大公会议关于神圣且裨益的复活节庆典之谕令者，若固执地坚持反对当时所正确定立的事，则应将他们绝罚并被逐出教会；此言针对平信徒。但若任何在教会中掌权者，无论主教、司铎或执事，在此谕令之后胆敢妄自判断，以致颠覆子民、扰乱教会，与犹太人同日守复活节，则神圣主教会议决定：此人此后应与教会隔绝，因为他不仅自招罪愆，且成为多人毁灭与颠覆之因；会议不仅免去其人本身的职务，且凡在其被免职后胆敢与之共融者，亦一并免职。被免职者甚至不得享有圣教规及天主的司祭职所共有的外在尊荣。\par

% structure: p0018 source=h3 target=subsection reason=relative-heading-rank; base=h2

\subsection{第二条}

凡进入天主的教会并聆听圣经，却不与众人一同祈祷，或由于某种紊乱而避开神圣的圣体圣事之领受者，应被逐出教会，直到他们告解、结出补赎之果、恳切祈求并获得赦免为止。与受绝罚者共融，或在私宅聚集与不在教会祈祷的人祈祷，或在另一教会中接纳不与另一教会聚会的人，皆为非法。若任何主教、司铎、执事或任何名列圣品者被发现与受绝罚者共融，亦应受绝罚，因其扰乱教会的秩序。\par

\EditorialNote{\Person{John Zonaras}: 在本条教规中，教父们指的是那些去教堂却不肯留下参加祈祷也不领受圣体的人，被某种乖戾或放任所阻止，也就是说，没有任何正当理由，而是轻率地，由于某种紊乱 [\GreekText{ἀταξίαν}]。}\par

% structure: p0021 source=h3 target=subsection reason=relative-heading-rank; base=h2

\subsection{第三条}

若有司铎、执事或任何属于司祭职的人，离开自己的堂区，出走并完全改变居所，在另一堂区长久停留，则不得再行圣职；尤其当他的主教召唤并催促他返回自己的堂区，而他却不服从时。若他坚持其紊乱行为，则应完全免除其职务，不再给他恢复的机会。若有另一位主教因此原因收留被免职者，则应由全体主教会议惩处，因其废除教会法规。\par

\EditorialNote{比较宗徒教规第15条和第16条。}\par

% structure: p0024 source=h3 target=subsection reason=relative-heading-rank; base=h2

\subsection{第四条}

若有主教被主教会议免职，或有司铎、执事被其主教免职，竟胆敢执行任何圣职，无论其按先前习惯负责主教、司铎或执事职务，则不得再有任何在另一主教会议上恢复的希望，亦无申辩之机；凡与他共融者，皆应被逐出教会，尤其若他们明知针对上述之人的判决，仍胆敢与之共融。\par

\EditorialNote{本条教规之所以引起主要兴趣，是因为它通常被认为是在反对\Person{圣亚大纳削}的党派鼓动下采纳的，后来又被用来反对\Person{圣金口若望}。但是，尽管这可能是某些人投票赞成、另一些人珍视它的秘密原因，但必须记住，其规定与宗徒教规相同，并且在迦克墩大公会议上作为第八十三条教规被宣读。\Person{雷米·塞利耶}（\WorkTitle{作者通史}，第659页）试图证明这不是\Person{圣金口若望}及其朋友所拒绝的教规，但\Person{赫费勒}认为他的立场完全站不住脚（\WorkTitle{大公会议史}，第二卷，第62页，注1），并引用了\Person{蒂勒蒙}（\WorkTitle{回忆录}，第329页，《论亚略派》，以及\Person{富克斯}的\WorkTitle{教会会议文献}，第二部分，第59页）。（比较宗徒教规第28条。）}\par

% structure: p0027 source=h3 target=subsection reason=relative-heading-rank; base=h2

\subsection{第五条}

若有司铎或执事藐视自己的主教，脱离教会，聚集私人集会，并设立祭台；且当主教召唤他，即使主教第一次、第二次召唤他，他仍拒绝被说服、不服从，则应完全免除其职务，不再有补救之方，亦不能恢复其品级。若他坚持扰乱并破坏教会，应由世俗权力以叛乱者身份予以纠正。\par

% structure: p0029 source=h3 target=subsection reason=relative-heading-rank; base=h2

\subsection{第六条}

若有人被自己的主教绝罚，则不得被他人接纳，直至他由其主教恢复，或直至主教会议召开时，他出庭申辩，并说服会议获得不同判决为止。此谕令适用于平信徒、司铎、执事以及所有登记在圣职人员名单中的人。\par

% structure: p0031 source=h3 target=subsection reason=relative-heading-rank; base=h2

\subsection{第七条}

没有和平信函者，不得接纳外人。\par

% structure: p0033 source=h3 target=subsection reason=relative-heading-rank; base=h2

\subsection{第八条}

乡村司铎不得颁发正典信函，或仅可将此类信函发送给邻近主教。但品行良好的乡村主教可以颁发和平信函。\par

这些\Quote{正典信函}在西方被称为\Quote{\OriginalTerm{格式信}{formatæ}}，在早期教会中将信徒联结在一起的巨大影响，最有力的证据莫过于\Person{背教者尤利安}曾试图在其帝国异教徒中引入类似制度。\par

\Quote{介绍信}（\OriginalTerm{介绍信}{\GreekText{ἐπιστολαὶ} \GreekText{συστατικαὶ}}）由圣保禄在\BibleRef{格后 3:1}中提及，读者可在\Person{J. J. 布伦特}的\WorkTitle{前三世纪的基督教会}（第二章）中找到关于此及类似主题的有趣评论。\par

凭借这些信函，即使平信徒也能在世界各地得到款待和照顾；多纳徒派被视为分裂者的标志之一，便是他们的正典信函仅在他们内部有效。\par

\Person{伪依西多尔}告知我们，在迦克墩大公会议上，君士坦丁堡主教\Person{阿提库斯}声称，在尼西亚大公会议上已同意所有此类信函应标有\OriginalTerm{圣父圣子圣神}{\GreekText{Π}. \GreekText{Υ}. \GreekText{Α}. \GreekText{Π}}（即圣父、圣子、圣神），并且据称（\Person{赫尔佐格}，\WorkTitle{真正百科全书}，词条\Quote{\OriginalTerm{格式信}{Literæ Formatæ}}）这一形式在六世纪的日耳曼文献中被发现。\par

如将在《加采东大公会议法令集》中所见，旧的名称\Quote{推荐信}仍在沿用，但本条法令及《劳迪西亚会议》第四十一条中使用了\Quote{正典信函}这一表述。至少在西方，这些信函加盖了教区主教印玺，以避免一切冒名顶替的可能。主任\Person{普林普特里}（我在此注释中严格遵循其见解）认为，使用教区印玺的最早证据见于\Person{奥古斯丁}（\WorkTitle{书信集}第五十九通，另编第二百一十七通）。他还引证了\Person{杜康热}，见其\WorkTitle{词汇}\OriginalTerm{格式}{Formatæ}条。\par

由于这些信函准许持信者参与共融，它们有时被称为\Quote{共融信函}（\OriginalTerm{共融信函}{\GreekText{κοινωνικαὶ}}），并由\Person{圣济利禄·亚历山大里亚}如此描述；还见于\Place{埃尔维拉}会议（第二十五条法令），以及\Person{圣奥古斯丁}（\WorkTitle{书信集}第四十三通，另编第一百六十二通）。\par

\Quote{和平信函}似乎具有施舍性质，持信者因此获得物质帮助。\Place{加采东}会议在其第十一条法令中规定，应将此类\Quote{和平信函}给予穷人，无论他们是圣职人员还是平信徒。本次会议的前一条法令也使用了相同的表述。\par

后来一种教会信函形式是我们所熟悉的\Quote{调职信函}。这一表述首次出现于\Place{特鲁洛}会议第十七条法令。关于此表述，\Person{斯维策}（\WorkTitle{词典}，见\OriginalTerm{调职}{\GreekText{ἀπολυτικὴ}}条）从上下文得出结论：\Quote{调职信函}仅用于永久变更教会住所，而\Quote{推荐信函}则给予那些暂时离开其教区的人。\par

% structure: p0043 source=h3 target=subsection reason=relative-heading-rank; base=h2

\subsection{第九条}

每位主教都应在每个省份中承认都主教，他须为整个省份操心；因为所有经商之人从各处汇集到省会。因此，规定他应享有优先地位，且其他主教未经他的同意，不得做任何非常之事（按照从我们教父时代起就通行的古老法令），除非是仅涉及他们自己特定堂区和所辖区域的事务。因为每位主教对自己堂区拥有权柄，既能以每人应尽的虔诚来管理它，又能为依赖其城市的所有区域提供照顾；可以按立司铎和执事；并以明智决定一切。但未经都主教同意，他不得进一步行事；都主教未经其他人同意亦然。\par

\EditorialNote{这里所提及的古老法令，几乎毫无疑问是指宗徒法令集第三十四条，因为在该条中读到了相同的规定（且几乎相同的措辞），只是此处阐述得稍为详尽；也没有任何其他法令在年代上早于此会议而包含这些规定，因此由此得出一个强有力的论据，支持宗徒法令集的完整性。}\par

比较宗徒法令第三十四条。\par

% structure: p0047 source=h3 target=subsection reason=relative-heading-rank; base=h2

\subsection{第十条}

神圣会议规定，在村庄和地区中的人，即那些被称为乡村主教者，即使他们已领受主教圣秩，也应遵守自己的辖区，管理辖下的教会，并满足于对这些教会的照管和治理；但他们可以按立读经员、副执事和驱魔员，并应满足于擢升这些人，但未经其所属城市主教（他和他的区域隶属于该主教）的同意，不得擅自按立司铎或执事。若他胆敢违反这些规定，他将被剥夺其所拥有的品级。乡村主教应由其所属城市的主教任命。\par

% structure: p0049 source=h3 target=subsection reason=relative-heading-rank; base=h2

\subsection{第十一条}

若有任何主教、司铎或任何属圣职人员者，未经省主教尤其是都主教的同意和信函，擅自前往皇帝处，此人应被公开撤职并驱逐，不仅逐出共融，亦逐出他所在的品级；因为他胆敢违反教会的法律，惊扰我们蒙天主宠爱的皇帝的圣听。但若有必须事务需要某人前往皇帝处，他应获得省主教和其他主教的建议和同意，并持他们的信函启程。\par

\EditorialNote{本条法令是早期教会为遏制后来所称的政府至上主义这一日益增长的倾向所做的重大努力之一。国家自成为基督徒后，不仅主动干预属灵事务，而且还出现了主教及其他圣职人员，他们无法以其他方式达到目的，便向世俗权力（通常是皇帝本人）上诉，从而威胁到整个教会的纪律，并废置了属灵会议的权威。从希腊注释家\Person{巴尔萨蒙}对本条法令的评注中，教会在这场斗争中屡遭挫折已显而易见。}\par

% structure: p0052 source=h3 target=subsection reason=relative-heading-rank; base=h2

\subsection{第十二条}

若有任何司铎或执事被其本主教撤职，或任何主教被会议撤职，胆敢惊扰皇帝的圣听，而他本应将案件提交更高级的主教会议，并请更多的主教审断他所认为合理的事，然后服从他们所进行的审查和决定；若他轻视这些，惊扰皇帝，他将不获任何宽恕，也不得有任何辩护的机会，亦无望将来复职。\par

\EditorialNote{通常认为，本条法令以及第四条、第十四条和第十五条是针对\Person{圣亚大纳修}的，并曾被\Person{圣金口若望}的敌人用来对付他。参见\Person{苏格拉底}\WorkTitle{教会史}第二卷第八章，以及\Person{索佐门}\WorkTitle{教会史}第三卷第五章；又同上书第七卷第二十章。}\par

% structure: p0055 source=h3 target=subsection reason=relative-heading-rank; base=h2

\subsection{第十三条}

任何主教不得擅自从一个省份前往另一个省份，并在教会中按立人员担任职务的尊位，即使有其他人与他同行，除非他应他前往地区的都主教和主教的书面邀请。但若他未经邀请，擅自进行任何按立，或处理与他无关的教会事务，他所行之事无效，他本人应因其违规和不合理的行为受到应有的惩罚，即被神圣会议立即撤职。\par

% structure: p0057 source=h3 target=subsection reason=relative-heading-rank; base=h2

\subsection{第十四条}

若一位主教因某些指控受审，且省主教们对他意见不一，有的认为被告无罪，有的认为有罪；为解决所有争议，神圣会议规定，都主教应从邻近省份另召若干主教，他们应加入审判并解决争议，从而与省主教们一同确认所决定的事。\par

% structure: p0059 source=h3 target=subsection reason=relative-heading-rank; base=h2

\subsection{第十五条}

若任何主教被指控，经省中所有主教审判，且所有主教一致对他做出相同裁决，则他不得再被他人审判，省主教们的一致判决应维持不变。\par

% structure: p0061 source=h3 target=subsection reason=relative-heading-rank; base=h2

\subsection{第十六条}

若任何无教区的主教，未经全体会议，擅自占据空缺的教会并夺取其宝座，即使所有被他篡夺管辖权的人民都选择他，他仍应被驱逐。全体会议应指都主教出席。\par

\EditorialNote{本条及下一条法令，由\Person{莱昂提乌斯主教}在\Place{加采东}大公会议上，从法令书中诵读；在该法令书中，按其所遵循的次序，本条称为第九十五条，下一条称为第九十六条。}\par

% structure: p0064 source=h3 target=subsection reason=relative-heading-rank; base=h2

\subsection{第十七条}

任何人若已接受主教圣秩，并被委任管理一个教区，却不愿接受其职务，并拒绝前往托付给他的教会，则应被绝罚，直至他受强制而接受为止，或直至省区全体主教会议对他做出决定为止。\par

% structure: p0066 source=h3 target=subsection reason=relative-heading-rank; base=h2

\subsection{第十八条}

如果有主教受祝圣于一个堂区，却未前往他被祝圣的堂区，并非由于他自身的过失，而是由于教众的拒绝或其他非出于他自身的原因，则应让他保有其品位和职务；只是他不得扰乱其所加入的教会的事务；并且他应遵守省区全体主教会议在审理案件后所做的任何决定。\par

% structure: p0068 source=h3 target=subsection reason=relative-heading-rank; base=h2

\subsection{第十九条}

\Emph{主教}不得在没有主教会议和省区都主教在场的情况下被祝圣。当都主教在场时，最好省区内所有在主内服务之同僚都与他一同聚集；都主教应以书信邀请他们。最好是全体集会；但如果这有困难，则不可或缺的是，多数应亲自出席或以书信参与选举，如此，任命应在多数在场或同意下进行；若违反这些规定而行，祝圣无效。若按所规定的法令做出任命，而有人出于好辩的天性提出反对，则多数的决定应生效。\par

\EditorialNote{选举主教的方法已规定于尼西亚法令第四条，但本条法令补充规定：若选举违反本决定之规定，则无效且无效；当选举者对所选举之人的意见分歧时，多数的票应生效。但当你听到本条法令说没有都主教在场就不应有选举时，你不可说他应亲身出席选举（因为这在其他法令中是被禁止的），而是说他在此出席是一种许可或劝说，没有它选举就不能进行。}\par

% structure: p0071 source=h3 target=subsection reason=relative-heading-rank; base=h2

\subsection{第二十条}

为了教会的益处和解决争论，法令规定：在每个省区内应每年举行两次主教会议（都主相应通知省区内主教），一次在复活节庆期第三周之后，以便会议在五旬节第四周结束；第二次在十月望日，即Hyperberetæus月的第十日（或第十五日），以便司铎、执事以及所有认为自己受到不公正对待的人可以求助于这些会议并获得会议的判决。但任何人无权在没有都主教座堂负责人在场的情况下自行举行会议。\par

\EditorialNote{尼西亚大公会议规定的在四旬期之前召开会议的时间在东方未被接受，主教们继续按照旧习惯在复活节后第四周举行会议，因为在四旬期之前的时间对于来自偏远城市的参会者来说常常带来极大的困难。}\par

% structure: p0074 source=h3 target=subsection reason=relative-heading-rank; base=h2

\subsection{第二十一条}

\Emph{主教}不得从一个堂区调往另一个堂区，无论是出于自己的建议擅自闯入，还是受教众强迫，或受主教们的强制；他应留在天主起初分配给他的教会，且不得从此调离，按照先前就此问题通过的法令。\par

% structure: p0076 source=h3 target=subsection reason=relative-heading-rank; base=h2

\subsection{第二十二条}

主教不得前往不属于自己管辖的他城，也不得进入不属于他的地区，无论是祝圣任何人，或任命司铎或执事到其他主教管辖范围内的地点，除非获得当地合法主教的同意。若有人妄图行此类事，祝圣无效，他本人应受会议惩罚。\par

% structure: p0078 source=h3 target=subsection reason=relative-heading-rank; base=h2

\subsection{第二十三条}

主教即使临终，也不得指定他人为自身继承人；若行此类事，指定无效。但应遵守教会法规：主教必须通过会议和主教们的裁决任命，主教们有权在安息于主怀中后擢升堪当之人。\par

\EditorialNote{本条文的规定极其重要。它显然旨在防止任何形式的裙带关系，并将空缺主教的任命完全置于都主教及其会议的自由选择之下。教会的历史及其现时的实践是对古代立法的奇特注释，而后来常见的\OriginalTerm{有继承权的助理主教}{coadjutor bishop cum jure successionis}似乎在某种程度上巧妙地规避了本条文。然而，凡·埃斯彭提醒读者注意希波的\Person{圣奥斯定}（他在第213封信中自述）的最有趣案例：他是如何被其前任选为希波主教的，而他和当时的主教都不知此事被法令禁止。以及当他年老时，人们希望他选一位主教帮助他直到去世并随后继任，他拒绝说：\Quote{在我自己身上应受责备的事，不可同样成为我儿子的污点。}他毫不迟疑地说出他认为最适合继任的人，但他补充说：\Quote{他应像现在这样作司铎，等天主旨意允许时，再作主教。}凡·埃斯彭补充说：这一切都应仔细阅读，从而知道圣奥斯定如何为主教和牧者们树立榜样，尽一切可能使他们在去世后，真正的牧人而非盗贼和狼进入他们的羊群，那些盗贼和狼会在短时间内摧毁他们用如此长久的辛劳所完成的一切。（参看\Person{欧瑟比}，\WorkTitle{教会史}第六卷第十一章及第三十二章。）}\par

% structure: p0081 source=h3 target=subsection reason=relative-heading-rank; base=h2

\subsection{第二十四条}

属于教会的一切，应当以良心和信德（信靠天主，他是万物的监察者和审判者）妥善保存给教会。这些事务应当在主教的判断和权下管理，主教被托付了整个教区和集会者的灵魂。但应当明确何为教会财产，使身边的司铎和执事知晓；以便他们确切知道哪些事物属于教会，不向他们隐瞒任何事，以便当主教离世时，属于教会的财产为人熟知，不致被侵吞或丢失，并使主教的私人财产不因被假托为教会财物而受干扰。因为主教遗赠其私人财产给所愿之人，而属于教会的则归教会保管，这是公正且为天主和人所喜悦的；如此，教会不受损失，主教不因教会利益之借口受损害，其亲属也不陷入诉讼，他本人在死后也不受责难。\par

% structure: p0083 source=h3 target=subsection reason=relative-heading-rank; base=h2

\subsection{第二十五条}

主教应有管理教会资金之权，以便以全备的虔诚和敬畏天主之心分施予所有需要者。若有需要，他可取用为自身必要所需及与他同住的弟兄们所需之物，使他们一无所缺，如同神圣宗徒所说，\BibleQuote{只要有吃有穿，就应当知足。}若他不以此满足，而将资金用于私用，且不经司铎和执事同意管理教会收入或农场租金，反而将权力交给自己的仆役、亲属、兄弟或儿子，以致教会账目暗中受损，则他本人应受省区会议的调查。反之，若主教或其司铎被控将属于教会之物（无论出自土地或其他教会资源）据为己有，以致穷人受压迫，并使账目和管理者招致控告和恶名，则他们也应受纠正，由神圣会议判定是非。\par

\EditorialNote{在拉贝版本的狄奥尼修斯译本中，本条末尾附加了以下文字：\Quote{参加本次会议的三十位主教签署。} 伊西多尔·梅卡托尔有更完整的文本，即：\Quote{我，欧瑟比，出席并签署本次神圣会议所制定的一切。 狄奥多勒、尼西塔、马其顿纽、阿纳托利、塔尔科迪曼图斯、埃塞雷乌斯、纳尔奇苏斯、欧斯塔基乌斯、赫西基乌斯、莫里基乌斯、保禄及其他三十位主教同意并签署。} 凡·埃斯彭指出这一附加未见希腊文本，亦未见马丁·布拉卡伦西斯本，并补充说这段文字不太可能与法令本身属同一时代。}\par

\SourceRecord{Translated by Henry Percival.；Nicene and Post-Nicene Fathers, Second Series；Vol. 14.；Edited by Philip Schaff and Henry Wace.；Buffalo, NY: Christian Literature Publishing Co.,；1900.；<http://www.newadvent.org/fathers/3805.htm>.}
